% Original analytic teaching diagram, 2026-11-09.
\definecolor{VB}{HTML}{7188A3}\definecolor{VBF}{HTML}{E5EAF0}
\definecolor{VG}{HTML}{829C8B}\definecolor{VGF}{HTML}{E8EEE9}
\definecolor{VY}{HTML}{B4A16F}\definecolor{VYF}{HTML}{F1ECDD}
\definecolor{VR}{HTML}{A7787D}\definecolor{VRF}{HTML}{F0E5E6}
\definecolor{VP}{HTML}{9588A5}\definecolor{VPF}{HTML}{ECE8F0}
\definecolor{VI}{HTML}{3B4252}\definecolor{VM}{HTML}{6B7280}\definecolor{VL}{HTML}{4C4D4F}
\begin{tikzpicture}[x=1mm,y=1mm,>=Stealth,font=\figurefont\fontsize{8}{10}\selectfont,
lab/.style={text=VM,font=\figurefont\fontsize{7.5}{9}\selectfont,align=center},
bx/.style={rounded corners=1.8pt,draw=VB,fill=white,line width=.75pt,minimum height=11mm,minimum width=30mm,align=center,inner sep=1.5mm},
key/.style={rounded corners=1.8pt,draw=VP,fill=VPF,line width=.75pt,minimum height=11mm,minimum width=32mm,align=center,inner sep=1.5mm},
ar/.style={-{Stealth[length=2mm]},draw=VL,line width=.8pt}]
\node[key] (Z) at (0,0) {可行候选对\\间隙为零};
\node[bx] (S) at (-38,-25) {Lagrange鞍点};
\node[bx] (K) at (38,-25) {KKT点};
\node[bx] (W) at (0,-25) {双方全局最优\\且零间隙};
\draw[ar] ([xshift=-2mm]Z.south)--node[lab,left]{弱对偶}([xshift=-2mm]W.north);
\draw[ar] ([xshift=2mm]W.north)--([xshift=2mm]Z.north|-Z.south);
\draw[ar] ([yshift=1.5mm]S.east)--([yshift=1.5mm]W.west);
\draw[ar] ([yshift=-1.5mm]W.west)--([yshift=-1.5mm]S.east);
\draw[ar,VL] ([yshift=1.5mm]K.west)--node[lab,above=.3mm]{凸}([yshift=1.5mm]W.east);
\draw[ar,VL] ([yshift=-1.5mm]W.east)--node[lab,below=.3mm]{可微}([yshift=-1.5mm]K.west);
\end{tikzpicture}
